
\section{CV-based QCPINN}
\label{app:appendix_a}

\begin{algorithm}[ht]
\footnotesize % or \small or \scriptsize
\caption{Continuous Variable Neural Network (CV-based QNN)}
\begin{algorithmic}[1]
\Require Qumodes: $num\_qumodes$, Layers: $num\_layers$, Device: $d$, Cutoff: $c$, Input: Tensor $x$
\Ensure Output: Tensor $y$

\State \textbf{Initialize Parameters:}
\State $\theta_1, \theta_2 \sim \mathcal{N}(0, 0.01\pi)$ (Interferometer), $r_{\text{disp}}, r_{\text{squeeze}} \sim \mathcal{N}(0, 0.001)$ (Nonlinear)
\State $\text{kerr} \sim \mathcal{N}(0, 0.001)$ (Kerr nonlinearity)

\State \textbf{Quantum Device:} Initialize $d$ with $num\_qumodes$, cutoff $c$.

\Procedure{QuantumCircuit}{$\text{input}$}
    \For{$i = 1$ to $num\_qumodes$}
        \State Encode Inputs with Displacements $D$: $D(\text{input}[i], 0)$ \Comment{only real amplitudes}
    \EndFor
    \For{$l = 1$ to $num\_layers$}
        \State Apply Interferometer($\theta_1[l]$)
        \State Apply Squeezing($r_{\text{squeeze}}[l], 0.0$)
        \State Apply Interferometer($\theta_2[l]$)
        \State Apply Displacement($r_{\text{disp}}[l], 0.0$)
        \State Apply Kerr($\text{kerr}[l] \times 0.001$)
    \EndFor
    \State Measure state $\langle \hat{q}_i \rangle$ for each wire.
    \State \Return Measurements
\EndProcedure

\Procedure{Interferometer}{$\theta$}
    \For{$l = 1$ to $num\_qumodes$}
        \For{$(q_1, q_2)$ in Pairs($num\_qumodes$)}
            \If{$(l + k) \% 2 \neq 1$}
                \State Apply Beamsplitter $B(\theta[n], 0.0)$
            \EndIf
        \EndFor
    \EndFor
    \For{$i = 1$ to $num\_qumodes - 1$}
        \State Apply Rotation $R_Z(\theta[i])$
    \EndFor
\EndProcedure

\Procedure{ForwardPass}{Inputs $x$}
    \State Initialize an empty list for outputs
    \For{each $s$ in $x$}
        \State $y \gets \text{QuantumCircuit}(s)$
        \State Append $y$ to outputs
    \EndFor
    \State \Return outputs
\EndProcedure

\end{algorithmic}
\label{alg:cv_qnn}
\end{algorithm}



